\section{Discussion}
\label{sec:discussion}

\subsection{One degree of freedom, not nine}

The three measurements in this paper are one mechanism observed three ways. AUC
says $p_{\mathrm{hit}}$ cannot rank moments within a video. Per-sample precision
being flat across the entire threshold grid, while recall falls monotonically, is
the same fact expressed in the metric we optimise: firing less often buys no
precision, which is exactly what a ranking statistic at chance predicts. And the
selection audit says that what remains --- a shallow, noisy dependence of F1 on
emission volume --- is not enough to identify nine separate operating points from
fifteen videos each.

What makes the result more than a null is that the same quantity \emph{is}
identifiable when the data are pooled. The single global threshold re-selects at
94\,\% against a 10\,\% chance rate, while the per-task fits sit at 25--61\,\%.
Nothing about the tasks changed between those two numbers; only the number of
videos backing each estimate did. This is a statistical power statement, not a
claim that the tasks are secretly identical --- their curves do differ, they are
simply each estimated too noisily to rank.

\subsection{Why a per-task table is the tempting wrong answer}

Per-task thresholds are attractive because the tasks genuinely differ in emission
rate, and because a per-task table always fits the fitting set better than a single
value. Our fitted arm does score highest in Table~\ref{tab:ablation}. It is only
when the comparison is made \emph{paired}, against the best single global
threshold, that the advantage turns out to be $+0.0075$ with an interval from
$-0.011$ to $+0.027$ --- and that comparison is in-sample, where nine free
parameters are given every advantage over one and still fail to separate. A study
that stopped at the four point estimates would have reported a per-task gain in
good faith.

\subsection{What would have to be true for per-task fitting to work}

The obstruction is ranking quality, not grid resolution, and this distinction has a
practical consequence: widening or refining the grid cannot help, because no
threshold recovers information the score does not carry. Two things would change
the answer. First, a gating score that actually ranks event-adjacent ticks above
quiet ones --- AUC meaningfully above $0.5$ --- would make precision respond to the
threshold, and per-task optima would become separable. Second, more videos per task
would shrink the per-task intervals; our own pooled result is the existence proof
that the estimator works given enough data. The two are not independent: the number
of videos needed scales with how weakly the score ranks.

\subsection{Implications for how streaming gates are reported}

We would draw a general methodological point from a specific negative result. A
per-task threshold table is a set of fitted parameters, and it should be held to
the standard applied to fitted parameters: a demonstration that the values recur
under resampling, and a comparison against the simplest model that could explain
the same data --- here, one global threshold. Both tests are cheap, since they
re-use the runs already performed and cost no additional GPU time. Neither is
standard practice, and in this study each of them independently overturns the
conclusion that the point estimates alone would have supported.

Finally, the two defects in Section~\ref{sec:method} deserve emphasis beyond this
system. The inert threshold would not have been visible in any aggregate metric:
the system ran, emitted, and scored plausibly, and the fitted values were simply
being ignored. It was found by checking that no fire had ever occurred below the
global constant, which is a check on the mechanism rather than on the result. A
grid search cannot detect that the parameter it is sweeping is disconnected.
